%%%%%%%%%%%%%%%%%%%%%%%%%%%%%%%%%%%%%%%%%%%%%%%%%%%%%%%%%
%%%%%%%%%%%%%%%%%%%         SOLUCION              %%%%%%%%%%%%%%%%%%%%%%%
%%%%%%%%%%%%%%%%%%%%%%%%%%%%%%%%%%%%%%%%%%%%%%%%%%%%%%%%%


\chapter{Discusión}
\label{chap:discussion}

\chapterintro{
En conjunto los tres artículos que conforman esta Tesis Doctoral estudian, tanto el potencial del microbioma como fuente de biomarcadores, como la aplicación de técnicas de ML para el análisis de datos del microbioma con el fin de apoyar procesos diagnósticos. Siguiendo el hilo argumental de los tres artículos y los resultados obtenidos, esta sección aborda aspectos específicos y transversales generados a lo largo de la investigación.
}



Esta Tesis Doctoral explora el microbioma como fuente de biomarcadores con potencial diagnóstico y emplea técnicas computacionales para convertir perfiles microbianos en señales útiles para la clasificación y la estratificación clínica. El uso de datos 16S rRNA se justifica por su accesibilidad, estandarización y amplia disponibilidad pública; cualidades que favorecen estudios a gran escala, validaciones externas y reproducibilidad, sin establecer una jerarquía frente a WGS, sino una elección pragmática alineada con objetivos y recursos. Bajo este marco, se integraron análisis de abundancia diferencial, selección de características y Machine Learning con un procesamiento explícitamente composicional, a fin de respetar la naturaleza relativa de los conteos y evitar pérdidas informativas por submuestreo.



En primera instancia, se decidió estudiar las enfermedades presentadas a lo largo de la Tesis Doctoral debido a su conocida y fuerte conexión con algún componente microbiano. En el caso de las enfermedades metabólicas, estas tienen una relación directa con el tracto intestinal y su correcto funcionamiento, un ambiente que presenta una gran carga microbiana. En cuanto a la VB, se sabe que la patogénesis de la enfermedad tiene un gran porcentaje microbiano. Al estudiar estas enfermedades \cite{paper1, paper2, paper3}, es fundamental considerar la naturaleza de las muestras utilizadas. Las muestras fecales, al ser deslocalizadas, reflejan una visión general del microbioma intestinal a lo largo de todo el tracto digestivo. Esto es crucial para enfermedades sistémicas y complejas como el SM o DT1/2, donde el impacto puede deberse a interacciones a gran escala dentro del sistema digestivo. En contraste, las muestras localizadas, como los hisopos vaginales, ofrecen un análisis más directo de un ecosistema microbiano específico y menos heterogéneo. Dado que el microbioma vaginal es menos diverso, las alteraciones en su composición son más fáciles de asociar directamente con alteraciones específicas en la salud vaginal. Este diseño permitió poner a prueba la versatilidad del marco metodológico en escenarios con distinta señal microbiana esperable: en VB se anticipaban patrones comunitarios estables y clínicamente relevantes; en DT1/DT2, firmas predictivas robustas, pero con relaciones más difusas, moduladas por múltiples covariables. Los resultados fueron coherentes con estas premisas.


Metodológicamente, se optó por ASVs para maximizar la resolución y reproducibilidad, y por DADA2 como método de inferencia por su modelado explícito de los errores provenientes de la secuenciación y su independencia del conjunto de muestras \cite{callahan2017exact, nearing2018denoising, caruso2019performance, prodan2020comparing}. Al mismo tiempo, la evidencia reunida en esta Tesis Doctoral muestra que OTUs y ASVs pueden converger en señales comunitarias comparables cuando el preprocesamiento es cuidadoso: en \cite{paper3}, un modelo de \textit{clustering} basado en OTUs se validó en datos procesados con ASVs, mostrando que los patrones comunitarios principales se mantienen entre ambos enfoques; y en \cite{paper2}, las OTUs son suficientes para extraer una firma predictiva con alto rendimiento. Esto está en línea con evidencias que indican que, pese a diferencias en riqueza o asignación taxonómica, la topología del \textit{clustering} y las señales globales tienden a ser comparables \cite{chiarello2022ranking}. Aunque la agrupación puede fusionar secuencias e incrementar la riqueza, las OTUs capturan patrones robustos a niveles superiores \cite{liu2023independent}. La elección, por tanto, debe alinearse con el objetivo: granularidad para detectar cambios sutiles (ASVs) o robustez comunitaria para clasificación (OTUs).


El procesamiento composicional fue clave; la mayor preocupación era mitigar la gran variabilidad en la profundidad de secuenciación entre muestras y representar fielmente las proporciones de los taxones microbianos. Frente a la rarefacción, la cual puede simplificar en exceso las comparaciones al eliminar datos potencialmente relevantes, transformaciones logarítmicas como CLR y métodos de abundancia diferencial compatibles (p. ej., LinDA) preservaron la información y aumentaron la sensibilidad \cite{gloor2016compositional, gloor2017microbiome, hong2022rarefy, zhou2022linda}. Gracias a ambos, en el primer trabajo \cite{paper1}, se extrajeron patrones significativos. En el caso de los otros trabajos \cite{paper2, paper3}, se comparó el rendimiento aplicando estas transformaciones frente a los datos crudos y se observó que el rendimiento se reducía al no aplicar dichas transformaciones. Además, en este último trabajo, aplicar estas transformaciones ayudó a reducir la variabilidad entre los diferentes enfoques utilizados (OTUs vs ASVs) a lo largo de las distintas cohortes. Aunque una parte de la literatura \cite{lambeth2015composition, dash2023functional, siptroth2023variation, chen2018vaginal, onywera2021predictive} sigue empleando normalizaciones que ignoran la composicionalidad, la evidencia acumulada en esta Tesis Doctoral respalda un tratamiento composicional consistente.



Con el objetivo de maximizar el espacio de descubrimiento, se adoptó un enfoque dirigido por datos en lugar de uno basado en conocimiento. El enfoque dirigido por conocimiento depende de hipótesis preexistentes, vías biológicas conocidas o bases de datos curadas. Aunque proporciona interpretabilidad biológica y es un excelente complemento para confirmar hallazgos, limita la capacidad de descubrimiento a lo ya conocido \cite{dakal2025advanced}. El enfoque dirigido por datos opera sin depender de suposiciones biológicas. Esto es crucial en la búsqueda exploratoria de nuevos biomarcadores, ya que permite analizar datos brutos y descubrir patrones ocultos o no lineales \cite{hemmati2024using}. Al elegir una metodología guiada por este tipo de enfoque, se seleccionó el ML, ya que se deseaba utilizar una metodología que manejase, en cierta medida, la alta variabilidad, dispersión y complejidad del microbioma, y que además tuviese explicabilidad. En el estado del arte actual, se utilizan muchos enfoques basados en el conocimiento en el campo del microbioma \cite{chong2017computational, bauer2018network, luo2024progress}. Sin embargo, a lo largo de los tres trabajos \cite{paper1, paper2, paper3}, al analizar datos brutos en lugar de un subconjunto basado en conocimiento previo, hemos descubierto tanto taxones ya mencionados en la literatura como otros poco documentados o desconocidos en la patología de estudio. Esto abre la puerta a investigar el papel de estos taxones en el desarrollo de enfermedades y su potencial como biomarcadores.


Para afrontar el desafío de la alta dimensionalidad en los datos del microbioma, se han implementado filtros rigurosos (conteos totales y prevalencia de taxones) al extraer la tabla de taxones y se han aplicado técnicas de selección de características que permiten reducir el número de taxones considerados en la construcción de modelos de ML, las cuales han demostrado ser efectivas en este contexto \cite{segata2011metagenomic}. Actualmente, el problema de la dimensionalidad está presente en el estudio de datos procedentes del microbioma, ya que numerosos artículos en el estado del arte realizan algún tipo de selección de características \cite{beck2015machine, biassoni2020gut, bakir2021discovering, bakir2022inflammatory, chavarria2025using}. Esta selección de características desempeña un papel importante en el segundo trabajo de esta Tesis Doctoral \cite{paper2}. Este artículo emplea la selección de características para identificar las especies más relevantes capaces de diferenciar las muestras entre sano y DT1. La estrategia seguida fue utilizar una técnica de selección de características de filtrado, lo que permite medir la relevancia de cada una de las especies con las clases de interés, identificando así cuántas y qué especies utilizar para construir el modelo predictivo. Los resultados muestran un aumento significativo en el rendimiento de los modelos al seleccionar un subgrupo de especies del número total disponible.


La variabilidad biológica y técnica se abordó mediante estrategias de validación estrictas. Por un lado, en el caso del aprendizaje no supervisado, el enfoque CCP ha permitido la construcción de subgrupos robustos y generalizables. En el aprendizaje supervisado, la implementación de una NCV garantiza un modelo robusto y generalizable, minimizando el riesgo de sobreajuste \cite{varma2006bias}. Esto se refleja en los trabajos \cite{paper2, paper3}, donde, gracias a esta validación, se ha manejado la diversidad en los datos de entrada, permitiendo la construcción de modelos con poca dispersión en su rendimiento, independientemente del conjunto de entradas. Esto se traduce en un modelo con altas capacidades para clasificar pacientes entre sanos y DT1 en el trabajo \cite{paper2}, y en un modelo en el trabajo \cite{paper3} que es capaz de predecir los VCS, incluso cuando los datos de entrada provienen de fuentes heterogéneas.

La elección algorítmica se adecuó al nicho y la tarea: modelos robustos a la dispersión y no linealidades (RF, XGBoost) frente a escenarios complejos, y modelos lineales (glmnet) cuando la estructura era más separable y se priorizaba la interpretabilidad \cite{goodswen2021machine, Lin2023MicrobiomeObesitySignatures}. El rendimiento superior de algoritmos no lineales como RF en el contexto de la predicción de DT1 \cite{paper2} se debe principalmente a la necesidad de modelar interacciones complejas del ecosistema microbiano intestinal, las cuales se ven agravadas por una alta variabilidad interindividual y las interacciones anfitrión-huésped. En contraste, el uso efectivo de un modelo lineal, como glmnet, en la estratificación de pacientes en función de los perfiles del microbioma vaginal \cite{paper3} refleja la naturaleza relativamente estructurada y linealmente separable de las comunidades microbianas vaginales, donde se pueden delimitar características clave con precisión. Las diferencias en el contexto de la enfermedad y la heterogeneidad de los nichos microbianos son, por lo tanto, fundamentales para explicar por qué los métodos no lineales son superiores en un escenario, mientras que los modelos lineales son más efectivos en otro.


Como se ha mencionado, el microbioma humano es un ecosistema altamente dinámico que exhibe variabilidad temporal, tanto a nivel composicional como funcional. Esta característica, reconocida en la literatura, plantea desafíos para la clasificación confiable de enfermedades utilizando datos del microbioma \cite{aldars2021systematic}. Sin embargo, los resultados obtenidos en esta Tesis Doctoral \cite{paper2, paper3}, al igual que otros en el estado del arte \cite{eck2017robust}, han demostrado que, incluso utilizando solo una muestra por paciente para la clasificación de enfermedades, se pueden lograr resultados robustos. Este hecho sugiere que la robustez en la clasificación entre estados patológicos no depende de capturar cada fluctuación transitoria, sino más bien de detectar los marcadores microbianos persistentes y dominantes que diferencian estos estados \cite{mehta2018stability}.


En \cite{paper1}, los pacientes con SM mostraron un aumento en la relación Firmicutes/Bacteroidetes, en línea con la hipótesis clásica sobre obesidad y fenotipos metabólicos \cite{santos2019role, baghi2022evaluation}. Se destacó una mayor prevalencia de \textit{Blautia}, cuya contribución a la disfunción metabólica podría implicar inflamación y disrupción de la barrera intestinal; además, en estudios previos, su abundancia se asoció con mayor severidad del SM en la enfermedad del hígado graso no alcohólico \cite{Liu2021, shen2017gut, Wang2016}. El enfoque seguido permitió la identificación de nuevos taxones que podrían desempeñar un papel protector en el contexto del SM. En particular, se relacionó la especie \textit{Parabacteroides merdae} con los controles sanos. Aunque no ha sido extensamente estudiada en humanos, trabajos recientes en modelos murinos sugieren que esta especie podría desempeñar un rol beneficioso a través del catabolismo de aminoácidos de cadena ramificada, cuya acumulación se ha relacionado con un aumento del riesgo metabólico \cite{qiao2022gut}. Además, estudios previos han reportado una reducción de ciertas especies del género \textit{Parabacteroides} en individuos con SM \cite{cui2022roles}.

En la progresión a DT2 se observó un aumento de la relación Bacteroidetes/Firmicutes y una disminución de Actinobacteria. Asimismo, se observa una pérdida de productores de butirato como \textit{Bifidobacterium} \cite{larsen2010gut, wu2010molecular, le2013alterations} y de taxones poco conocidos en este contexto: \textit{Lachnospiraceae GCA-900066575} y \textit{Defluviitaleaceae UCG-011} (potenciales fermentadores de fibra y productores de AGCCs). Estudios muy recientes realizados en ratones señalan que \textit{GCA-900066575} se relaciona negativamente con niveles de glucosa en sangre en ayunas y la resistencia a la insulina \cite{liu2023vitamin}. Simultáneamente, se observó un incremento en los taxones asociados con la patogénesis, como \textit{Dorea}, \textit{Prevotella} y \textit{Dialister invisus}. En el caso específico de \textit{Dorea longicatena}, se ha encontrado que esta especie está enriquecida en pacientes con DT2. Se ha observado una correlación clínica significativa entre su presencia y el desarrollo de la enfermedad \citep{qin2014alterations}. En cuanto a \textit{Prevotella}, varios estudios han establecido una asociación entre un aumento de esta bacteria y la intolerancia a la glucosa, así como la diabetes gestacional \citep{zhang2013human, egshatyan2016gut, hasain2020gut}. Por otro lado, la especie \textit{Dialister invisus} se ha vinculado con un aumento de la permeabilidad intestinal y la diabetes \citep{maffeis2016association, zheng2018gut}.



El segundo trabajo \cite{paper2} muestra que la profundidad taxonómica del estudio es un factor determinante para capturar la señal asociada a la T1D. Se observa que el rendimiento de los tres algoritmos utilizados mejora a través de seis niveles taxonómicos, desde filo hasta especie, con mejoras del 30\% entre ambos niveles. Esto sugiere que la asociación con la enfermedad se basa en la dinámica o interacción entre especies, en lugar de en desequilibrios de niveles superiores. La firma metagenómica identificada se alinea con asociaciones previamente descritas a nivel de especie. \textit{Bacteroides dorei} se ha vinculado con la seroconversión y con etapas preclínicas de autoinmunidad \cite{davis2014bacteroides}. \textit{Prevotella copri} y \textit{Bacteroides vulgatus} se han asociado con perfiles metabólicos de resistencia a la insulina \cite{pedersen2016human}. En cuanto a \textit{Bifidobacterium animalis}, se han observado efectos favorables en ensayos clínicos (principalmente en DT2) sobre control glucémico y parámetros cardiometabólicos \cite{stenman2016probiotic, khiaolaongam2025bifidobacterium}. En el eje intestino-inmunidad, \textit{Bacteroides thetaiotaomicron} se ha asociado con fenotipos de enfermedad celíaca y presenta una capacidad inmunomoduladora relevante \cite{sanchez2010intestinal}. Por último, la dinámica de \textit{Escherichia coli} productora de amiloide se ha relacionado con vías prodiabéticas y con la seroconversión en cohortes pediátricas \cite{tetz2019type}. Los modelos RF entrenados con esta firma presentan un gran rendimiento (0,94 de \textit{accuracy}) tanto al clasificar entre individuos sanos y pacientes de DT1, como al añadir una clase intermedia de seroconvertidos. Esto sugiere una fuerte asociación entre estas especies bacterianas y la susceptibilidad a DT1.

En este contexto, la mayoría de los estudios describen alteraciones composicionales y funcionales asociadas a la seroconversión y DT1 \cite{vatanen2018human, kostic2015dynamics}. De forma similar, otros estudios poblacionales \cite{belteky2023infant} han aplicado algoritmos de ML (RF) para la clasificación de taxones en búsqueda de biomarcadores; sin embargo, no se detalla su rendimiento predictivo para poder comparar con este trabajo. Desde la publicación de este segundo trabajo, en los últimos años se han publicado diferentes estudios con nuevas cohortes de infantes (como PRJNA875929 \cite{belteky2023infant} y PRJNA702261 \cite{del2022gut}), dando pie a realizar estudios de metacohortes con enfoques similares al presentado en esta Tesis Doctoral, como es el caso de \cite{liu2024predicting}.



En el tercer trabajo \cite{paper3}, los subgrupos obtenidos mediante \textit{VIBES} (VCS I–IV) mostraron una alta estabilidad en la cohorte de descubrimiento, con un ARI medio de 0,914. Esto sugiere que la estratificación en cuatro grupos está determinada por la estructura intrínseca de los datos, y no por el tamaño muestral. Además, la distribución de los subtipos fue consistente en cinco cohortes externas. En la validación externa del modelo, la BA fue superior a 0.91 en todas las cohortes, y el \textit{kappa} fue mayor a 0,83.

Desde la perspectiva biológica, las betas del modelo cuantifican la contribución relativa de cada especie en el clasificador para determinar la pertenencia a un subtipo. VCS-I, dominado por \textit{Lactobacillus crispatus}, se asocia con eubiosis: producción elevada de ácido láctico, pH bajo y exclusión competitiva de patobiontes \cite{ravel2011vaginal, srinivasan2012bacterial}. VCS-II, con \textit{Lactobacillus iners} (junto con \textit{Aerococcus christensenii}, \textit{Campylobacter ureolyticus}, \textit{Staphylococcus haemolyticus}), sugiere un estado de transición o vulnerabilidad aumentada \cite{macklaim2011crossroads, holm2023lactobacillus}. En VCS-III, las betas relevantes incluyen \textit{Lactobacillus iners}, \textit{Prevotella bivia}, \textit{Campylobacter ureolyticus} y \textit{Finegoldia magna}. Las especies del género \textit{Prevotella} pueden elevar el pH mediante la producción de aminas biógenas (putrescina, cadaverina, trimetilamina) y amonio, favoreciendo el crecimiento de \textit{Gardnerella} y promoviendo la transición hacia la disbiosis \cite{miller2016lactobacilli, yeoman2013multi}. Finalmente, VCS-IV es el subtipo más heterogéneo, evidenciando altas betas para \textit{Atopobium vaginae}, \textit{Sneathia sanguinegens}, \textit{Mycoplasma hominis} y \textit{Campylobacter ureolyticus}, y en menor medida para \textit{Lactobacillus gasseri}. La especie \textit{Atopobium vaginae} está asociada a \textit{biofilms}, un elevado pH y recurrencias debido a su menor sensibilidad al metronidazol \cite{swidsinski2008adherent, bradshaw2006association}; mientras que \textit{Sneathia} expresa factores de virulencia con potencial inflamatorio \cite{fredricks2005molecular, harwich2012genomic}.

Metodológicamente, VIBES y VALENCIA ofrecen visiones complementarias. \textit{VALENCIA} utiliza \textit{k-nearest neighbors} y hasta 199 taxones (clase–especie) para clasificar en 5 CSTs/subclases \cite{france2020valencia}; VIBES emplea un conjunto de 20 especies para 4 subtipos. Las diferencias en algoritmo, espacio de características y estructura de clases permiten captar señales no redundantes: perfiles sanos/asintomáticos de VALENCIA tienden a mapear con VCS-I, mientras que CSTs como III, IV-A o V no se alinean de forma directa con VCS-II/III, aportando granularidad complementaria \cite{ravel2011vaginal}. En la cohorte longitudinal tratada con metronidazol, el modelo combinado (VALENCIA + VIBES) mejoró la predicción de respuesta frente a cada enfoque por separado, y los no respondedores se concentraron en VCS-IV (y en menor medida en VCS-III), lo que sugiere que los subtipos de VIBES añaden información clínica independiente y potencialmente accionable para la estratificación basada en microbioma.



Las principales limitaciones se relacionan con el tamaño y la distribución muestral en \cite{paper1}, la ausencia de validación pediátrica independiente equivalente en \cite{paper2} (debido a la falta de cohortes comparables) y la diversidad geográfica y el panel de especies compartidas en \cite{paper3}. Más allá de las particularidades de cada estudio, la investigación del microbioma enfrenta fuentes de variabilidad técnica. Estas fuentes surgen a lo largo del flujo de trabajo: la elección de la región del gen 16S rRNA, los kits de extracción de ADN, las tecnologías de secuenciación empleadas (Illumina, PacBio, Ion Torrent), la estrategia de agrupación seguida (OTUs frente a ASVs), el método bioinformático para realizar dicha agrupación y las bases de datos de referencia utilizadas para la asignación taxonómica.

En resumen, las principales contribuciones de la presente Tesis Doctoral se pueden sintetizar en cuatro ejes fundamentales: la capacidad del microbioma como fuente de biomarcadores, el gran potencial del uso de técnicas de ML en el análisis del microbioma, la identificación de posibles biomarcadores microbianos con potencial de aplicación clínica en el contexto de la DT1/DT2 y el SM, y el desarrollo de una nueva estrategia de estratificación del microbioma vaginal en el contexto de la VB. Estas aportaciones evidencian la utilidad práctica del enfoque propuesto, instando a continuar explorando el microbioma como una herramienta clave para la personalización de intervenciones médicas. Este compendio, en conjunto, sienta las bases para un futuro en el cual las tecnologías emergentes y la integración de datos complejos configuren el panorama de la investigación biomédica, aportando beneficios tangibles a la salud pública y a la práctica clínica moderna. En el futuro, es crucial avanzar hacia la integración de plataformas analíticas para unificar datos, combinando información del microbioma y del anfitrión para entender mejor las interacciones biológicas y sus implicaciones en la salud. Además, investigar el microbioma como barrera en la eficacia de tratamientos y analizar la resistencia antimicrobiana facilitarán el desarrollo de estrategias innovadoras.